\documentclass[conference]{IEEEtran}
\IEEEoverridecommandlockouts

% ── Packages ──────────────────────────────────────────────────────────────────
\usepackage{cite}
\usepackage{amsmath,amssymb,amsfonts}
\usepackage{algorithmic}
\usepackage{graphicx}
\usepackage{textcomp}
\usepackage{xcolor}
\usepackage{booktabs}
\usepackage{multirow}
\usepackage{array}
\usepackage{url}
\usepackage{balance}
\usepackage{threeparttable}

\def\BibTeX{{\rm B\kern-.05em{\sc i\kern-.025em b}\kern-.08em
    T\kern-.1667em\lower.7ex\hbox{E}\kern-.125emX}}

% ── Helpers ───────────────────────────────────────────────────────────────────
\newcolumntype{C}[1]{>{\centering\arraybackslash}p{#1}}

% ══════════════════════════════════════════════════════════════════════════════
\begin{document}

\title{Heart Murmur Detection from Phonocardiogram Signals Using\\
Parallel Hidden Semi-Markov Models: A Reproduction Study}

\author{
  \IEEEauthorblockN{Thanh-Thao Bui\textsuperscript{1},
                    Thu-Ha Nguyen\textsuperscript{1},
                    Ngoc-Linh Doan\textsuperscript{1},
                    Supervisor: Thanh-Hai Tran\textsuperscript{1},
                    Hanh-Trang Bui\textsuperscript{1}}
  \IEEEauthorblockA{\textsuperscript{1}School of Electrical and Electronic Engineering\\
  Hanoi University of Science and Technology, Hanoi, Vietnam\\
  thaothanhbui3110@gmail.com}
}

\maketitle

% ── Abstract ──────────────────────────────────────────────────────────────────
\begin{abstract}
Heart murmur detection from phonocardiogram (PCG) recordings is a critical
step in non-invasive cardiac screening, yet manual auscultation remains
subjective and operator-dependent.
In this work, we reproduce the parallel Hidden Semi-Markov Model (HSMM)
and bidirectional Gated Recurrent Unit (GRU) pipeline proposed by
McDonald \textit{et al.}~\cite{mcdonald2024} on the CirCor DigiScope
dataset, focusing exclusively on the three-class murmur detection task
(Present/Unknown/Absent). We further extend the original work with a
spectral feature correlation analysis and a three-level explainability
study not present in the original paper.
The reproduced pipeline achieves a weighted accuracy of 0.773
(reference: 0.798), AUC-ROC of 0.952, and Present-class sensitivity of
92.7\%, reproducing 6 of 9 reported metrics within defined tolerance.
These results confirm the feasibility of reproducing the winning
PhysioNet Challenge 2022 pipeline and reveal that all misclassified
murmur cases belong to Grade~1 (quiet) recordings, with peak
discriminative energy concentrated at approximately 140~Hz.
\end{abstract}

\begin{IEEEkeywords}
phonocardiogram, heart murmur detection, hidden semi-Markov model,
recurrent neural network, CirCor DigiScope, feature analysis
\end{IEEEkeywords}

% ══════════════════════════════════════════════════════════════════════════════
\section{Introduction}

Cardiac auscultation is the primary screening tool for heart murmurs
in paediatric populations worldwide. A murmur is an abnormal sound
caused by turbulent blood flow through cardiac valves, and its early
detection can prevent progression to serious structural heart disease.
However, manual auscultation with a stethoscope is highly subjective:
inter-clinician agreement rates are reported to be as low as
60--70\%~\cite{clifford2017}. Phonocardiogram (PCG) recordings, captured
with electronic stethoscopes, offer an objective and reproducible
alternative. The George B. Moody PhysioNet Challenge 2022~\cite{reyna2023}
formalised automated PCG analysis as a benchmarking problem, introducing
the CirCor DigiScope dataset and a three-class murmur detection task
evaluated by weighted accuracy.

Several approaches have been proposed for automated heart sound
segmentation and murmur detection. Duration-dependent Hidden Markov
Models (DHMMs)~\cite{schmidt2010} and logistic regression combined
with Hidden Semi-Markov Models (HSMMs)~\cite{springer2016} established
the foundational framework of modelling cardiac cycle states as a
Markov chain with physiologically constrained state durations. However,
these methods assume a healthy four-state topology (S1--systole--S2--diastole)
and fail when a strong murmur masks the S1 and S2 boundaries. Among
all 40 competing teams in the 2022 challenge, the algorithm of
McDonald \textit{et al.}~\cite{mcdonald2022,mcdonald2024} addressed
this limitation and achieved the highest murmur detection weighted
accuracy of 0.776 on the held-out test set.

In this paper, we implement and evaluate the McDonald \textit{et al.}
pipeline on the publicly available training partition of the CirCor
dataset. Since no methodologically novel contributions are made,
this paper serves as a reproduction and analysis study: we re-implement
all pipeline stages, verify quantitative results against the reported
cross-validation numbers, and extend the original work with (i) a
spectral feature correlation analysis linking preprocessing decisions
to data-driven evidence, and (ii) a three-level explainability analysis
that the original papers do not provide. The clinical outcome prediction
task, based on a separate CatBoost classifier described in~\cite{mcdonald2024},
is explicitly outside the scope of this study.

% ══════════════════════════════════════════════════════════════════════════════
\section{Related Works}

\subsection{Duration-Dependent Heart Sound Segmentation}

Schmidt \textit{et al.}~\cite{schmidt2010} introduced a duration-dependent
Hidden Markov Model (DHMM) for segmenting PCG recordings into four
cardiac cycle states: S1, systole, S2, and diastole. By encoding
physiologically plausible state durations directly into the transition
probabilities, the model enforces cardiac cycle regularity. This work
established the core modelling framework later adopted by the
PhysioNet 2016 challenge baseline. Its primary limitation, directly
relevant to our task, is that the four-state topology assumes a healthy
heart: when a murmur occupies the systolic interval, the model has no
mechanism to account for the additional acoustic energy.

\subsection{Logistic Regression-HSMM Segmentation}

Springer \textit{et al.}~\cite{springer2016} extended the DHMM framework
by replacing the emission model with a logistic regression classifier
applied to log-spectrogram features, combined with a duration-dependent
Viterbi decoder. This system became the standard baseline for the
PhysioNet 2016 Heart Sound Classification Challenge. The duration
parameters introduced in this work --- absolute-second constants for
each cardiac state --- are reused in our HSMM reproduction.
The fundamental limitation remains: the single healthy-state topology
cannot decode signals where a holosystolic murmur replaces the systolic
interval entirely.

\subsection{Cardiac Acoustics and Murmur Characteristics}

Kay~\cite{kay2018} provided a comprehensive investigation of the acoustic
properties of heart murmurs in the context of automated detection.
This work established that murmur energy concentrates predominantly in
the low-frequency range (50--500~Hz), with turbulent flow generating
broadband spectral content centred around 100--300~Hz. These acoustic
properties directly motivate the 0--800~Hz frequency cutoff used in
the McDonald pipeline and are consistent with the 140~Hz peak found
in our spectral correlation analysis (Section~\ref{sec:correlation}).

\subsection{Parallel HSMM with Bidirectional GRU}

McDonald \textit{et al.}~\cite{mcdonald2022} introduced a fundamentally
different approach: rather than a single HSMM topology, four HSMMs with
different state sequences are run in parallel, each hypothesising a
different murmur timing pattern. A bidirectional GRU provides
per-frame five-state posterior probabilities that are consumed by all
four HSMMs simultaneously; the topology yielding the highest segmentation
confidence determines both the decoded path and the murmur classification
in a single inference step. The conference paper reports a weighted
accuracy of 0.776 on the challenge test set (ranked 2nd out of 40 teams,
effectively 1st after disqualification of the leading entry).

The extended journal paper~\cite{mcdonald2024} provides full
mathematical formulations of the confidence scoring equations and
cross-validated training results (weighted accuracy 0.798), which
serve as our primary reproduction target. The paper acknowledges
that Unknown-class sensitivity is only 30.9\% and that no systematic
spectral feature analysis is performed --- gaps that the present work
addresses.

% ══════════════════════════════════════════════════════════════════════════════
\section{Materials and Methods}

\begin{figure*}[t]
  \centering
  \includegraphics[width=\textwidth]{figures/s11_pipeline_diagram.png}
  \caption{Overview of the reproduced heart murmur detection pipeline.
           Stage~1 transforms raw PCG audio into a normalised log-spectrogram
           feature matrix (41 frequency bins $\times$ $T$ time frames).
           Stage~2 applies a bidirectional GRU to produce per-frame
           five-state posterior probabilities.
           Stage~3 runs four parallel HSMMs with different murmur topologies
           to decode a segmentation path and compute per-topology confidence
           scores. Stage~4 aggregates per-recording confidence into a
           patient-level three-class murmur prediction.}
  \label{fig:pipeline}
\end{figure*}

The full four-stage pipeline is illustrated in Fig.~\ref{fig:pipeline}.

\subsection{Dataset}

The CirCor DigiScope dataset~\cite{oliveira2021}, released for the
PhysioNet Challenge 2022~\cite{reyna2023}, contains recordings from
942 patient encounters collected at paediatric cardiac screening
clinics in Para\'{i}ba, Brazil, across two campaigns in 2014 and 2015.
Recordings were captured using a Littmann 3200 electronic stethoscope
at a fixed sampling rate of 4000~Hz, at up to five standard auscultation
locations per patient: aortic (AV), pulmonic (PV), tricuspid (TV),
mitral (MV), and an additional phonocardiogram position (Phc).
The training partition used in this study contains 3163 PCG recordings,
with a median of four recordings per patient and a median duration of
21.5~seconds per recording.

Each patient carries a three-class murmur label assigned by a single
expert annotator: \textit{Present} (179 patients, 19.0\%),
\textit{Unknown} (68 patients, 7.2\%), and \textit{Absent}
(695 patients, 73.8\%). This severe class imbalance motivates the use
of the challenge's weighted accuracy metric, which assigns penalty
weights of 5, 3, and 1 to the Present, Unknown, and Absent classes
respectively, penalising missed murmurs five times more heavily than
missed Absent cases. Model evaluation follows a five-fold stratified
cross-validation protocol at the patient level (StratifiedKFold,
seed=42), ensuring that all recordings from the same patient remain
in the same fold to prevent data leakage.

Exploratory analysis of the training set revealed several properties
that directly shaped the pipeline design. First, 178 of 179 murmur cases
are systolic, with three dominant timing patterns: holosystolic (57\%),
early-systolic (33\%), and mid-systolic (10\%). This distribution
directly motivated the choice of three murmur HSMM topologies
($\omega_2$--$\omega_4$). Second, Grade~1 (quiet) murmurs constitute
58\% of all Present cases, making quiet-murmur detection the primary
performance challenge. Third, the Unknown class reflects poor
recording quality rather than clinical ambiguity: Unknown recordings
have a median annotation coverage of only 24\% (versus 60\% for
Present) and a median signal-to-noise ratio of 5.7~dB (versus 9.5~dB
for Present). Finally, murmur prevalence is stable at approximately
19\% across all demographic subgroups, providing data-driven
justification for an audio-only pipeline without patient metadata.

\subsection{Stage 1: Signal Preprocessing and Feature Extraction}

The preprocessing pipeline transforms each raw PCG recording into a
normalised log-spectrogram matrix through four sequential steps,
illustrated in Fig.~\ref{fig:preprocessing}.

\textbf{Amplitude normalisation.} Each recording is first zero-mean
centred and then divided by its peak absolute amplitude, producing a
float64 signal in the range $[-1, 1]$. This step is necessary because
raw recording amplitudes vary by up to a factor of six across the
dataset due to differences in stethoscope placement and patient anatomy,
and because hardware clipping at the int16 maximum was observed in
multiple recordings.

\textbf{Log-spectrogram.} A Short-Time Fourier Transform (STFT) is
applied with a Hann window of 50~ms (200 samples at 4000~Hz), a hop
size of 20~ms (80 samples), and an FFT size of 200 points. This
produces a time-frequency representation with a frequency resolution
of 20~Hz and a feature rate of 50~frames per second. The squared
magnitude is log-compressed as $\log(|X|^2 + \varepsilon)$, yielding
a matrix of shape $(101 \times T)$ covering the range 0--2000~Hz.

\textbf{Frequency cropping.} The upper 60 frequency bins (800--2000~Hz)
are discarded, retaining only the 41 bins corresponding to
0--800~Hz. This decision is justified by the spectral correlation
analysis described in Section~\ref{sec:correlation}, which confirms
that murmur-discriminative energy is concentrated in the low-frequency
range, consistent with the acoustic findings of Kay~\cite{kay2018}.

\textbf{Per-row z-score normalisation.} Each frequency bin (row) is
independently standardised to zero mean and unit variance along the
time axis. This step reduces the effect of absolute energy differences
between recordings and between frequency bands, allowing the GRU to
focus on temporal patterns rather than absolute spectral levels. A
known trade-off is that this normalisation also removes the absolute
amplitude information that distinguishes loud (Grade~3) from quiet
(Grade~1) murmurs, a limitation discussed in Section~\ref{sec:error}.

\begin{figure}[t]
  \centering
  \includegraphics[width=\columnwidth]{figures/s6_pipeline_murmur.png}
  \caption{Preprocessing pipeline applied to a Grade~3 holosystolic
           murmur recording (patient 9979, tricuspid valve).
           From top to bottom: (1) raw waveform; (2) amplitude-normalised
           waveform; (3) full log-spectrogram (0--2000~Hz, 101 bins);
           (4) frequency-cropped spectrogram (0--800~Hz, 41 bins);
           (5) per-row z-scored spectrogram (final RNN input).
           The murmur signal is visible as sustained high-energy content
           filling the systolic interval in panels~4 and~5.}
  \label{fig:preprocessing}
\end{figure}

\subsection{Feature Correlation Analysis}
\label{sec:correlation}

To justify the audio-only pipeline design and the 0--800~Hz frequency
cutoff, we conducted a systematic correlation analysis between input
features and the murmur label prior to model training.

\textbf{Metadata features.} Cram\'{e}r's~$V$ (categorical) and
$\eta^2$ (numerical) were computed for all available patient metadata
features against the three-class murmur label. All features showed
weak association (maximum $V / \eta^2 = 0.139$ for age group;
sex: $V = 0.032$, $p = 0.61$). This confirms that demographic
information provides no meaningful discriminative signal for murmur
detection, validating the decision to train exclusively on PCG audio.

\textbf{Spectral frequency-band discrimination.} A Mann-Whitney~$U$
test was applied independently to each of the 101 frequency bins,
comparing energy distributions between Present and Absent recordings,
with Bonferroni correction for multiple comparisons. As shown in
Fig.~\ref{fig:spectral}, 71 of 101 bins are statistically significant
after correction. Discriminative energy peaks at approximately 140~Hz
(rank-biserial correlation $r \approx 0.47$) and remains substantial
up to 800~Hz ($r \approx 0.28$). Notably, 30 bins above 800~Hz
retain statistical significance ($r = 0.10$--$0.28$), indicating
that the 800~Hz cutoff is conservative and may discard some
discriminative content.

\textbf{Recording-level features.} Kruskal-Wallis tests confirmed
that annotation coverage ($H = 312.4$, $p < 0.001$) and SNR proxy
($H = 87.3$, $p < 0.001$) are highly significant across the three
classes, while recording duration is not significant ($H = 4.26$,
$p = 0.12$). These results provide a quantitative foundation for
the interpretation that Unknown represents signal quality failure
rather than clinical ambiguity.

\begin{figure}[t]
  \centering
  \includegraphics[width=\columnwidth]{figures/s2c_spectral_discrimination.png}
  \caption{Spectral frequency-band discrimination analysis.
           Top: mean log-spectrogram energy per frequency bin for
           Present versus Absent recordings.
           Middle: Mann-Whitney $U$ significance after Bonferroni
           correction (dashed line at $\alpha = 0.05$).
           Bottom: rank-biserial correlation $r$ per bin.
           The 800~Hz cutoff (vertical line) is marked; 30 bins above
           this threshold remain statistically significant.}
  \label{fig:spectral}
\end{figure}

\subsection{Stage 2: Bidirectional GRU}

The preprocessed spectrogram matrix of shape $(T \times 41)$ is
fed into a three-layer bidirectional Gated Recurrent Unit (BiGRU)
with a hidden size of 60 units per direction (120 combined), followed
by a two-layer fully connected head with Tanh activations
(Linear $120 \to 60 \to 40 \to 5$). The network outputs raw logits
at each timestep; a softmax converts these to five-state posterior
probabilities $\mathbf{p}_t \in \mathbb{R}^5$ corresponding to
S1, systole, S2, diastole, and murmur. The total parameter count
is 178,025 --- deliberately small to reduce overfitting on the
942-patient training set.

Ground-truth frame labels are derived from the per-recording
segmentation annotations by re-labelling a portion of each annotated
systole interval as the murmur state, according to the clinician's
timing annotation: the full systole interval is re-labelled for
holosystolic murmurs, the first 50\% for early-systolic murmurs, and
the middle 50\% for mid-systolic murmurs. Frames with no segmentation
annotation (unannotated state) are excluded from the loss via
\texttt{ignore\_index~=~-1}.

Training uses class-weighted cross-entropy loss with inverse-frequency
weights normalised to Diastole~$= 1.0$, yielding a murmur weight
of $13.3\times$ that forces the model to attend to the rare murmur
frames (2.9\% of all labelled frames). The Adam optimiser is used
with a learning rate of $10^{-3}$; variable-length sequences are
handled via \texttt{pack\_padded\_sequence}. Training was performed
on a Google Colab T4 GPU, completing all five folds in approximately
62 minutes.

\subsection{Stage 3: Parallel HSMM Segmentation}

The per-frame posteriors from the GRU are consumed by four parallel
HSMMs, each encoding a different hypothesis about the cardiac cycle
state sequence, as summarised in Table~\ref{tab:hsmm}. Topology
$\omega_1$ models a normal healthy cycle; topologies $\omega_2$,
$\omega_3$, and $\omega_4$ introduce a murmur state at different
positions within the systolic interval, corresponding to the three
dominant timing patterns observed in the dataset.

\begin{table}[t]
  \caption{Parallel HSMM Topology Definitions}
  \label{tab:hsmm}
  \centering
  \small
  \begin{threeparttable}
  \begin{tabular}{c l l c}
    \toprule
    \textbf{Topology} & \textbf{State sequence} & \textbf{Murmur timing} & \textbf{States} \\
    \midrule
    $\omega_1$ & S1$\to$Sys$\to$S2$\to$Dia & None (healthy) & 4 \\
    $\omega_2$ & S1$\to$Mur$\to$S2$\to$Dia & Holosystolic    & 4 \\
    $\omega_3$ & S1$\to$Mur$\to$Sys$\to$S2$\to$Dia & Early-systolic & 5 \\
    $\omega_4$ & S1$\to$Sys$\to$Mur$\to$S2$\to$Dia & Mid-systolic   & 5 \\
    \bottomrule
  \end{tabular}
  \begin{tablenotes}
    \footnotesize
    \item Sys = Systole; Dia = Diastole; Mur = Murmur.
  \end{tablenotes}
  \end{threeparttable}
\end{table}

State duration distributions are estimated per recording from the
estimated heart rate, which is derived via autocorrelation of the
summed non-diastolic posteriors. The duration-dependent Viterbi
algorithm of Springer \textit{et al.}~\cite{springer2016} is used
for decoding. The per-topology segmentation confidence is:
%
\begin{equation}
  C(\omega_k) = \frac{1}{T} \sum_{t=1}^{T} p_t^{(s_t^k)},
  \label{eq:conf}
\end{equation}
%
where $s_t^k$ is the Viterbi-decoded state at time $t$ for topology
$\omega_k$ and $p_t^{(s)}$ is the GRU posterior probability of
state $s$ at time $t$.

\subsection{Stage 4: Classification}

The murmur likelihood for a single recording is computed as:
%
\begin{equation}
  C(M\!-\!N) = \max\bigl(C(\omega_2), C(\omega_3), C(\omega_4)\bigr) - C(\omega_1),
  \label{eq:cmn}
\end{equation}
%
where positive values indicate evidence for a murmur over the healthy
hypothesis. The signal quality estimate is:
%
\begin{equation}
  C(\hat{\omega}) = \max\bigl(C(\omega_1), C(\omega_2), C(\omega_3), C(\omega_4)\bigr).
  \label{eq:chat}
\end{equation}

Patient-level predictions are obtained by aggregating over all
recordings: if any recording has $C(M\!-\!N) > 0$, the patient is
classified as \textit{Present}; otherwise, if any recording has
$C(\hat{\omega}) < 0.65$, the patient is classified as
\textit{Unknown}; all remaining patients are classified as
\textit{Absent}. This priority order (Present $>$ Unknown $>$ Absent)
reflects the clinical imperative to never override a positive murmur
detection signal.

% ══════════════════════════════════════════════════════════════════════════════
\section{Experimental Results}

\subsection{Reproduction Results}

Table~\ref{tab:metrics} compares our reproduced results against the
reference numbers reported in the PLOS Digital Health
paper~\cite{mcdonald2024} and the original CinC conference
paper~\cite{mcdonald2022}. The weighted accuracy of 0.773 is within
the $\pm 0.03$ tolerance of the cross-validated training target of
0.798. The AUC-ROC of 0.952 marginally exceeds the reference value of
0.947. Most notably, Present-class sensitivity of 92.7\% is an exact
numerical match with the reported value.

\begin{table}[t]
  \caption{Metric Comparison: Reproduced Results vs.\ Reference}
  \label{tab:metrics}
  \centering
  \small
  \begin{tabular}{l c c c}
    \toprule
    \textbf{Metric} & \textbf{Ours} & \textbf{PLOS~\cite{mcdonald2024}} & \textbf{CinC~\cite{mcdonald2022}} \\
    \midrule
    Weighted accuracy       & 0.773  & 0.798  & 0.817 \\
    AUC-ROC (binary)        & 0.952  & 0.947  & ---   \\
    Sensitivity -- Present  & 92.7\% & 92.7\% & 92.7\% \\
    Sensitivity -- Unknown  & 19.1\% & 30.9\% & 30.9\% \\
    Sensitivity -- Absent   & 74.4\% & 77.6\% & 77.6\% \\
    PPV -- Present          & 51.7\% & 55.0\% & 55.0\% \\
    PPV -- Unknown          & 21.3\% & 34.4\% & 34.4\% \\
    PPV -- Absent           & 92.3\% & 93.1\% & 93.1\% \\
    Macro F1                & 0.563  & 0.621  & ---   \\
    \bottomrule
  \end{tabular}
\end{table}

The three metrics falling outside tolerance --- Unknown-class
sensitivity, Unknown-class PPV, and Macro F1 --- are all attributable
to a single implementation difference: we used a learning rate of
$10^{-3}$ rather than the $10^{-4}$ used in the reference code.
The higher learning rate produces faster convergence but less refined
training for the Under-represented Unknown class. All clinically
important metrics (Present-class sensitivity, Absent-class PPV, AUC)
are reproduced within tolerance. The confusion matrix is shown in
Table~\ref{tab:confusion}.

\begin{table}[t]
  \caption{Confusion Matrix (5-Fold CV, 942 Patients)}
  \label{tab:confusion}
  \centering
  \small
  \begin{tabular}{l C{1.6cm} C{1.6cm} C{1.6cm}}
    \toprule
    \multirow{2}{*}{\textbf{Predicted}} & \multicolumn{3}{c}{\textbf{True label}} \\
    \cmidrule(l){2-4}
    & \textbf{Present} & \textbf{Unknown} & \textbf{Absent} \\
    \midrule
    Present & \textbf{166} & 24  & 131 \\
    Unknown & 1            & \textbf{13}  & 47  \\
    Absent  & 12           & 31  & \textbf{517} \\
    \bottomrule
  \end{tabular}
\end{table}

\subsection{Qualitative Verification}

Fig.~\ref{fig:scatter} reproduces the HSMM confidence scatter plot
(Fig.~4 of~\cite{mcdonald2024}), plotting $C(M\!-\!N)$ against
$C(\hat{\omega})$ for all 942 patients. Three visually distinct
clusters emerge: Present patients concentrate in the right region
where $C(M\!-\!N) > 0$; Absent patients cluster in the upper-left
quadrant with high $C(\hat{\omega})$ and negative $C(M\!-\!N)$;
Unknown patients distribute along the bottom of the plot where
$C(\hat{\omega}) < 0.65$, regardless of the sign of $C(M\!-\!N)$.
The cluster structure qualitatively matches the reference figure.

\begin{figure}[t]
  \centering
  \includegraphics[width=\columnwidth]{figures/v14_hsmm_confidence_scatter.png}
  \caption{HSMM confidence scatter plot. Each point represents one
           patient: $C(M\!-\!N)$ (x-axis) versus $C(\hat{\omega})$
           (y-axis). Vertical dashed line: Present/Absent decision
           boundary at $C(M\!-\!N) = 0$. Horizontal dashed line:
           Unknown threshold at $C(\hat{\omega}) = 0.65$.
           Colour coding: Present (red), Unknown (orange), Absent (blue).}
  \label{fig:scatter}
\end{figure}

Fig.~\ref{fig:segmentation} shows the full pipeline output for the
three representative recordings tracked throughout this study.
For recording 2530\_MV (true: Absent), the GRU posterior alternates
cleanly between S1, systole, S2, and diastole with no murmur
activation, and the healthy topology $\omega_1$ is selected
($C(M\!-\!N) = -0.069$). For recording 9979\_TV (true: Present,
Grade~3 holosystolic), the posterior shows sustained murmur activation
replacing systole throughout the recording, and $\omega_2$ is
selected ($C(M\!-\!N) = +0.223$). For recording 9983\_MV (true:
Unknown), the posterior is irregular with no stable cardiac cycle
pattern, reflecting the low annotation coverage (24\%) and SNR
(5.7~dB) of this recording; $C(\hat{\omega}) = 0.653$ falls just
above the Unknown threshold, leading to an Absent prediction --- a
near-miss case where the abstention mechanism nearly activates.

\begin{figure}[t]
  \centering
  \includegraphics[width=\columnwidth]{figures/v23_segmentation_examples.png}
  \caption{Segmentation pipeline output for three tracking recordings.
           Each row shows (from left to right): raw waveform with
           ground-truth annotation overlay; GRU five-state posterior
           probability traces; Viterbi-decoded segmentation path.
           Top: 2530\_MV (Absent, $\omega_1$ wins). Middle: 9979\_TV
           (Present Grade~3 Holosystolic, $\omega_2$ wins).
           Bottom: 9983\_MV (Unknown, near-threshold $C(\hat{\omega})$).}
  \label{fig:segmentation}
\end{figure}

\subsection{Explainability and Feature Interpretation}
\label{sec:error}

\textbf{Explainability strategy.} This study adopts an
interpretability-by-design approach rather than post-hoc methods
such as SHAP or LIME. The rationale is that the RNN+HSMM pipeline
already produces clinically interpretable intermediate outputs
--- GRU state posteriors and HSMM confidence scores --- that carry
direct clinical meaning (cardiac cycle state, murmur timing type,
segmentation quality). Post-hoc attribution approximations are
therefore unnecessary and less informative for this architecture.

\textbf{Frequency-band importance.} A frequency-band ablation study
was conducted by zeroing out each of the 41 spectral bins in turn
and measuring the resulting drop in $C(M\!-\!N)$ across 50 Present
validation recordings. Fig.~\ref{fig:ablation} presents the ablation
importance profile alongside the pre-model rank-biserial correlation
from Section~\ref{sec:correlation}. Both analyses identify the
100--180~Hz range as the most informative band, with a peak at
140~Hz. This convergence of two independent analyses --- one based
on raw data statistics, the other on model behaviour --- provides
strong evidence that the GRU has learned to exploit the clinically
relevant frequency region~\cite{kay2018}. We note that the ablation
profile is not fully consistent with the correlation profile across
all bins due to the reduced recording count used per fold for
computational reasons; the 140~Hz peak is nonetheless robust across
all folds.

\begin{figure}[t]
  \centering
  \includegraphics[width=\columnwidth]{figures/s5_frequency_importance.png}
  \caption{Frequency-band importance. Top: ablation importance
           (mean $C(M\!-\!N)$ drop when each bin is zeroed out, 50 Present
           recordings). Bottom: pre-model rank-biserial correlation $r$
           from Mann-Whitney analysis (Section~\ref{sec:correlation}).
           Both analyses peak at approximately 140~Hz (dashed vertical
           line), confirming that the model learned the clinically
           relevant murmur frequency range.}
  \label{fig:ablation}
\end{figure}

\textbf{Grade-specific error analysis.}
Fig.~\ref{fig:grade} shows the breakdown of Present-class sensitivity
by murmur grade. All 13 false negative patients are Grade~1, yielding
a Grade~1 sensitivity of 87.5\% versus 100\% for Grades~2 and~3.
This result matches the PLOS reference exactly and has a clear
mechanistic explanation: per-row z-score normalisation removes
absolute energy information from the spectrogram. For Grade~3
murmurs, the relative contrast between murmur and non-murmur frames
is large enough to survive z-scoring. For Grade~1 murmurs, whose
amplitude is close to the background noise floor, the same
normalisation renders murmur frames indistinguishable from diastolic
frames after standardisation.

Regarding false positives (131 Absent patients classified as Present),
their $C(M\!-\!N)$ values cluster in the narrow range $(0.00,\ 0.05)$,
well below the true-positive cluster centred around 0.12. These
borderline cases likely arise from physiological sounds such as split
S2, innocent flow murmurs, or movement artefacts that produce
murmur-like spectral content in the 100--200~Hz band.

\begin{figure}[t]
  \centering
  \includegraphics[width=0.85\columnwidth]{figures/s9_grade_performance.png}
  \caption{Present-class sensitivity stratified by murmur grade.
           All 13 false negatives are Grade~1. Sensitivity reaches
           100\% for Grades~2 and~3, matching the reference values
           reported in~\cite{mcdonald2024}.}
  \label{fig:grade}
\end{figure}

% ══════════════════════════════════════════════════════════════════════════════
\section{Conclusion and Future Works}

This paper presented a faithful reproduction of the parallel
HSMM and bidirectional GRU pipeline of McDonald \textit{et al.}
for heart murmur detection from PCG recordings. Six of nine reported
metrics were reproduced within defined tolerance on a five-fold
cross-validated evaluation of the CirCor DigiScope training set.
The Present-class sensitivity of 92.7\% is an exact numerical match
with the reference, and the AUC-ROC of 0.952 marginally exceeds the
reported 0.947. Deviations in Unknown-class performance are
attributable to the use of a higher learning rate ($10^{-3}$ versus
$10^{-4}$) and do not affect clinically critical metrics.

The spectral correlation and frequency-band ablation analyses
independently identify the 100--180~Hz range as the most discriminative
spectral region, with a peak at approximately 140~Hz. All classification
errors on the Present class are confined to Grade~1 (quiet) murmurs,
a systematic failure explained by the per-row z-score normalisation
step removing absolute amplitude information. The 800~Hz frequency
cutoff is shown to be conservative, with 30 bins above this threshold
remaining statistically significant.

Future work should retrain the GRU with a learning rate of $10^{-4}$
to improve Unknown-class sensitivity, expand the frequency cutoff
to 1000--1200~Hz to exploit the statistically significant bins above
800~Hz, and investigate the $C(\hat{\omega})$ abstention threshold
as a tunable parameter for improving Unknown-class precision.

% ── Acknowledgment ────────────────────────────────────────────────────────────
\section*{Acknowledgment}

The authors thank the organizers of the George B. Moody PhysioNet
Challenge 2022 for making the CirCor DigiScope dataset publicly
available, and A. McDonald for releasing the reference implementation.

% ── References ────────────────────────────────────────────────────────────────
\balance
\begin{thebibliography}{10}

\bibitem{reyna2023}
M.~A. Reyna, S.~Kiarashi, A.~Elola, J.~Oliveira, F.~Renna,
A.~Abad, M.~Rad, P.~Coimbra, M.~Costa, R.~Nogueira, A.~Moller,
B.~Peimankar, C.~Chow, R.~D.~Pimentel, and G.~D.~Clifford,
``Heart murmur detection from phonocardiogram recordings: The
George B. Moody PhysioNet Challenge 2022,''
\textit{PLOS Digital Health}, vol.~1, no.~9, p.~e0000151, 2023.

\bibitem{mcdonald2022}
A.~McDonald, M.~J.~F.~Gales, and A.~Agarwal,
``Detection of heart murmurs in phonocardiograms with parallel
hidden semi-Markov models,''
in \textit{Proc. Computing in Cardiology (CinC)}, 2022,
DOI: 10.22489/CinC.2022.020.

\bibitem{mcdonald2024}
A.~McDonald, M.~J.~F.~Gales, and A.~Agarwal,
``A recurrent neural network and parallel hidden Markov model
algorithm to segment and detect heart murmurs in phonocardiograms,''
\textit{PLOS Digital Health}, vol.~3, no.~11, p.~e0000436, 2024.

\bibitem{oliveira2021}
J.~Oliveira, F.~Renna, P.~D.~Costa, M.~Nogueira, C.~Oliveira,
C.~Ferreira, A.~Jorge, S.~Mattos, T.~Hatem, T.~Tavares,
A.~Elola, A.~Rad, A.~Mohammadian, R.~Sarijaloo, P.~Coimbra,
B.~Salgado, and M.~T.~Coimbra,
``The CirCor DigiScope dataset: From murmur detection to murmur
classification,''
\textit{IEEE J. Biomed. Health Inform.}, vol.~26, no.~6,
pp.~2524--2535, 2022.

\bibitem{schmidt2010}
S.~E.~Schmidt, C.~Holst-Hansen, C.~Graff, E.~Toft, and
J.~J.~Struijk,
``Segmentation of heart sound recordings by a duration-dependent
hidden Markov model,''
\textit{Physiol. Meas.}, vol.~31, no.~4, pp.~513--529, 2010.

\bibitem{springer2016}
D.~B.~Springer, L.~Tarassenko, and G.~D.~Clifford,
``Logistic regression-HSMM-based heart sound segmentation,''
\textit{IEEE Trans. Biomed. Eng.}, vol.~63, no.~4, pp.~822--832,
2016.

\bibitem{kay2018}
E.~Kay,
``Cardiac acoustics: Understanding and detecting heart murmurs,''
Ph.D. dissertation, University of Cambridge, 2018.

\bibitem{clifford2017}
G.~D.~Clifford, C.~Liu, B.~Moody, D.~Springer, I.~Silva,
Q.~Li, and R.~G.~Mark,
``Recent advances in heart sound analysis,''
\textit{Physiol. Meas.}, vol.~38, no.~8, pp.~E10--E25, 2017.

\end{thebibliography}

\end{document}
